\section{Transmission through a cell}
\label{sec:transmission}

Throughout this section, the geometry and branch parameters are those of
\Cref{cell:geometry,cell:profile,cell:branch}. Constants may depend on the
fixed energy, disorder, profile, corner-box exponents, and parameters
already chosen, but not on exterior matrices or source vectors. A
polynomial bound means $C(1+L)^C$ with this convention; different occurrences
may have different exponents. All assertions concern sufficiently large
$L$ after the parameters have been fixed.

\subsection{Exterior data and the two branch estimates}

Suppose that the root is the only cell port in $V\cup B$, and that the
source port lies outside this set. Artificial probes with a prescribed
mask may be present. Eliminate the exterior physical sites. Its effect
on $B$ is a fixed complex-symmetric passive matrix $T$, and the source
gives a fixed vector $k$, including its input multiplier. The remaining
boundary matrix is
\[
 \mathcal T=T+\diag_B\bigl(v_i-i\tau a_i(v_i)\bigr).
\]
The exterior is held fixed throughout the cell integration, as are the
four corner boxes. In particular, no bound on $\norm{T}$ or $\norm{k}$ is
assumed.

Let $P$ also denote its range, equipped with a real orthonormal basis.
Eliminating the four missing boundary directions gives
\begin{equation}
 \widehat T=(P\mathcal T^{-1}P)^{-1},\qquad
 \widehat k=\widehat T P\mathcal T^{-1}k,\qquad
 k_*=(\widehat T-iI_P)^{-1}\widehat k,\qquad
 \rho=\norm{k_*}^2.
 \label{trans:normalization}
\end{equation}
These formulas are used at their almost-everywhere inverses. The
missing-direction block is generically invertible by
\Cref{imp:schur-interface}. Schur complementation
preserves symmetry and passivity, and gives the identity
\[
 k_*=P(\mathcal T-iP)^{-1}k.
\]
Thus $\rho$ is exactly the weight in \Cref{imp:weighted}, and
\begin{equation}
 \norm{(\widehat T-iI_P)^{-1}}\le1.
 \label{trans:absorbed-exterior}
\end{equation}
Hereafter $M$ is restricted to $P$, and is extended by a zero row at the
root when it is regarded as a map into $V$.

Let $G$ be the inverse on $V\oplus P$ with the root open, reference
potential $c(v)$, and root strength $\sigma=\epsilon^2$. Its internal
artificial absorptions are constant during the branch integration:
the true values lie on the plateau specified in \Cref{cell:branch}.
Set
\begin{equation}
 u=\sqrt\sigma\,Ge_0,\qquad g=\sigma G_{00},\qquad
 t=\widehat k^\top u_P.
 \label{trans:amplitude}
\end{equation}
Since $k$ includes the source input multiplier, the root-to-source
scattering amplitude satisfies $\abs{S_{\mathrm{source},a}}^2=4\abs t^2$.
The sign convention for the source does not affect this identity.
The inverse exists almost everywhere even though $c(v)$ depends on
other potentials. Indeed,
absorption forces a kernel vector to vanish at the root. Restoring its
true real diagonal therefore preserves that vector. After eliminating
the fixed boundary data, the real part is a fixed matrix plus the true
real diagonals on $V$; their absolutely continuous joint law makes its
determinant nonzero almost everywhere. This also justifies the
almost-everywhere slices used below.

\begin{proposition}[Branch transmission]\label{trans:cell}
Choose the branch parameters in the order described below and fix
$c_*>2\nu$. For $0\le\tau\le e^{-c_*L}$, uniformly in the fixed exterior
data and allowed corner boxes,
\begin{align}
 \E\abs{t}^2
 &\ge e^{-C(\log L)^2}\,\E_B\rho,
 \label{trans:lower}\\
 \E\bigl[\abs\mu\,\abs{t}^2\bigr]
 &\le e^{-cL}\,\E\abs{t}^2.
 \label{trans:relative}
\end{align}
Here $\mu$ is the conditional single-root closure coefficient of
\Cref{mat:closure,mat:cutoff}. The expectation is the branch expectation,
and $\E_B$ averages its boundary law. The hidden root value has already
been averaged in $\mu$; averaging it again does not change any other
quantity in these formulas.
\end{proposition}

\subsection{The cutoff error and the small-response region}

We first derive \Cref{trans:relative} assuming \Cref{trans:lower}.
Fix the boundary data and all jitters, and vary the common real shift
$r$ on $Z^+$. Write $P_+$ for the coordinate projection onto $Z^+$,
extended by zero on the root and on $P$. Since
$c(v)=E+p^\top A^{-1}p$ and $A'=P_+$ on $J$, differentiation gives
\begin{equation}
 c'=-\norm{P_+A^{-1}p}^2,\qquad
 g'=-u^\top P_+u-\sigma^{-1}c'g^2.
 \label{trans:shift-derivatives}
\end{equation}
The profile estimate in \Cref{cell:residual}, together with
$\abs{\psi_i}\le C\epsilon$ on $Z^+$, implies
$\abs{c'}/\sigma\le C(1+L)^C$. Root absorption and
\Cref{mat:ward} imply
\begin{equation}
 \abs g^2\le\operatorname{Im}g\le1.
 \label{trans:g-ward}
\end{equation}

Let $I_r=[-L^{-D_1},L^{-D_1}]$. The derivative of a continuous
determinant phase of $G^{-1}$ is
\[
 \Phi'(r)=\tr(P_+\operatorname{Im}G P_+)
       +\frac{c'}{\sigma}\operatorname{Im}g.
\]
The phase has range at most $\pi\dim(V\oplus P)$ after strict
regularization. Integrating, using \Cref{trans:g-ward}, and then
removing the regularization as in \Cref{mat:phase} gives
\begin{equation}
 \int_{I_r}\tr(P_+\operatorname{Im}G P_+)\,\dd r
 \le C(1+L)^C.
 \label{trans:phase-trace}
\end{equation}
In particular, Ward and symmetry bound the integral of
$\norm{P_+u}^2$ by the same polynomial. We need an additional factor
when $g$ is small.

Put $D=-\operatorname{Im}G^{-1}$. Complex symmetry makes $D$ real
symmetric and gives
\[
 u-\overline u=2iG^*Du,\qquad
 u^*Du=\operatorname{Im}g.
\]
For each coordinate, Cauchy--Schwarz in the form $D$ bounds the square
of the corresponding entry of $G^*Du$ by
$(\operatorname{Im}G)_{ii}\operatorname{Im}g$. Consequently, with
entrywise imaginary part on vectors,
\begin{equation}
 \norm{P_+\operatorname{Im}u}^2
 \le (\operatorname{Im}g)
       \tr(P_+\operatorname{Im}G P_+).
 \label{trans:imaginary-vector}
\end{equation}
Take $q=\sqrt\sigma=\epsilon$, and set
$\Omega_q=\{r\in I_r:\abs{g(r)}\le q\}$, excluding null singular
sets. \Cref{trans:phase-trace,trans:imaginary-vector} give
\[
 \int_{\Omega_q}\norm{P_+\operatorname{Im}u}^2\,\dd r
 \le C(1+L)^Cq.
\]

There are only polynomially many component intervals of $\Omega_q$.
To see this, $A$ is affine in $r$, $c$ is a rational function obtained
by inverting $A$, and $G$ is a further matrix inverse in cell dimension
$O(L^2)$. Thus numerator and denominator degrees of $g$ are bounded
by a fixed polynomial in $L$. The inequality $\abs g^2\le q^2$
becomes a real polynomial inequality after clearing the squared
denominator. Its intervals and the poles are therefore polynomially
many; the values of the exterior coefficients do not affect this
degree bound. Integrate the real part of
\Cref{trans:shift-derivatives} on these intervals, first truncating
away from any endpoints where necessary. Every endpoint value of
$\operatorname{Re}g$ has modulus at most $q$. Hence
\[
 \int_{\Omega_q}\operatorname{Re}(u^\top P_+u)\,\dd r
 \le C(1+L)^Cq+C(1+L)^Cq^2\abs{I_r}.
\]
Using
$\norm{P_+u}^2=\operatorname{Re}(u^\top P_+u)
 +2\norm{P_+\operatorname{Im}u}^2$, we obtain
\begin{equation}
 \int_{\Omega_q}\norm{P_+u}^2\,\dd r
 \le C(1+L)^Cq.
 \label{trans:small-response}
\end{equation}
Normalized uniform integration in $r$ only introduces the polynomial
factor $(2L^{-D_1})^{-1}$.

The $Z$ rows of the equation $G^{-1}u=\sqrt\sigma e_0$ involve only
$u_P$ and $P_+u$. Their coefficients are uniformly bounded, and $M$
is bounded below on $P$. The $P$ rows then give, respectively,
\[
 \norm{u_P}\le C\norm{P_+u},\qquad
 \widehat T u_P=-M^\top u_J,\qquad
 \norm{\widehat T u_P}\le C\norm{P_+u}.
\]
It follows from symmetry and \Cref{trans:normalization} that
\begin{equation}
 \abs t
 =\abs{k_*^\top(\widehat T-iI_P)u_P}
 \le C\sqrt\rho\,\norm{P_+u}.
 \label{trans:source-control}
\end{equation}
The cutoff bounds in \Cref{mat:cutoff} now yield
\begin{equation}
 \E\bigl[\abs\mu\,\abs t^2\bigr]
 \le Cq\E\abs t^2+C(1+L)^Cq\E_B\rho.
 \label{trans:relative-prelower}
\end{equation}
The first term treats $\abs g>q$; the second follows from
\Cref{trans:small-response,trans:source-control} on $\Omega_q$.
Since $q=e^{-\nu L}$, substitution of \Cref{trans:lower} absorbs
the factor $e^{C(\log L)^2}$ and proves \Cref{trans:relative} for
some fixed $c>0$.

\subsection{Regular boundary and shift data}

We turn to \Cref{trans:lower}. By \Cref{imp:weighted}, it is enough
to work on boundary data satisfying
\begin{equation}
 \norm{\widehat T^{-1}}\le C(1+L)^C,
 \label{trans:good-boundary}
\end{equation}
since this event carries a fixed positive fraction of $\E_B\rho$.
Fix such data, assume $\rho>0$, and initially set every jitter to zero.
Define the boundary Schur matrix
\[
 \mathcal B=\widehat T-M^\top A^{-1}M.
\]
We claim that at least half of the uniform $r$-interval satisfies
\begin{equation}
 \norm{\mathcal B^{-1}}\le C(1+L)^C.
 \label{trans:regular-r}
\end{equation}
All exponents in this assertion can be fixed before $D_2,D_3$.

Indeed,
\[
 \mathcal B'=M^\top A^{-1}P_+A^{-1}M\ge cI_P.
\]
For the last inequality, put $w=A^{-1}Mx$. The $Z$ rows of $Aw=Mx$
have bounded coefficients, use only $P_+w$, and give
$\norm{Mx}\le C\norm{P_+w}$. The lower singular-value bound for
$M$ gives $\norm{P_+w}\ge c\norm x$, which proves the assertion.
For $\xi>0$ let $K_\xi=(\mathcal B-i\xi I_P)^{-1}$. If
$\norm{\mathcal B^{-1}}>\xi^{-1}$, inverse difference implies
$\norm{K_\xi}\ge(2\xi)^{-1}$. Ward then implies
\begin{equation}
 \xi\tr\operatorname{Im}K_\xi\ge\frac14.
 \label{trans:regular-test}
\end{equation}
The same conclusion follows at a singular matrix from a null vector.
The determinant phase derivative of $\mathcal B-i\xi I_P$ is at least
$c\tr\operatorname{Im}K_\xi$. Thus phase integration bounds the
Lebesgue measure of the failed event in \Cref{trans:regular-r} by
$C\xi\dim P$. Dividing by $2L^{-D_1}$ and choosing $\xi$ to be a
sufficiently small inverse polynomial proves the claim. The only
dependence on the exterior here is through coefficients of a passive
matrix; its norm never enters.

For regular zero-jitter data, omit both the root and artificial
absorption in $J$, and write
\begin{equation}
 R_0=\begin{pmatrix}A&M\\M^\top&\widehat T\end{pmatrix}^{-1},
 \qquad \Gamma=(R_0)_{JJ},\qquad
 z_0=(R_0)_{JP}\widehat k,\qquad
 U_0=\Gamma p/\epsilon.
 \label{trans:rootless-fields}
\end{equation}
The subscript $0$ denotes the absence of internal artificial
absorption; the boundary matrix $\widehat T$ is unchanged. These
definitions will also be used after the jitters have been varied.
The field $z_0$ is the response in $J$ to the exterior source. Define
its normalized coupling to the root by
\[
 F_0=p^\top z_0/\epsilon.
\]
If we adjoin the calibrated open root, still omitting the artificial
absorptions on $V$, block inversion gives the analogue of $t$ for this
auxiliary system:
\[
 -\frac{F_0}{(c(v)-E-i\sigma-p^\top\Gamma p)/\sigma}.
\]
The prescribed internal absorptions will be restored below.
We first obtain a lower bound for this numerator using
the shell jitters, then control the denominator. A diagonal change at
a site $j$ changes $F_0$ through the product $(U_0)_j(z_0)_j$.
We will find a source component near the boundary and a large component
of $U_0$ at logarithmic depth, and use the jitters to propagate the
source component to that deeper site.

Block inversion through $A$ and $\mathcal B$ gives
polynomial bounds on $R_0$ and $\Gamma$. Moreover,
\begin{equation}
 \Gamma^{-1}=A-M\widehat T^{-1}M^\top
 \label{trans:gamma-inverse}
\end{equation}
has polynomial norm by \Cref{trans:good-boundary}.

The source normalization requires some care because $\widehat k$
and $\widehat T$ need not be bounded. Put
$C_{\rm ext}=\widehat T-iI_P$, $\Delta=\mathcal B-C_{\rm ext}
=iI_P-M^\top A^{-1}M$, and $b_0=\mathcal B^{-1}\widehat k$.
The identities
\begin{equation}
 b_0=(I_P-\mathcal B^{-1}\Delta)k_*,\qquad
 k_*=(I_P+C_{\rm ext}^{-1}\Delta)b_0
 \label{trans:source-comparison}
\end{equation}
use only polynomially bounded $\Delta,\mathcal B^{-1}$ and the
norm bound \cref{trans:absorbed-exterior}. Consequently,
\begin{equation}
 (1+L)^{-C}\sqrt\rho\le\norm{b_0}
 \le C(1+L)^C\sqrt\rho,\qquad
 \norm{z_0}\le C(1+L)^C\sqrt\rho.
 \label{trans:source-size}
\end{equation}
Since $Az_0=-Mb_0$, the sparse $Z$ rows also imply
$\norm{P_+z_0}\ge c\norm{b_0}$. There is therefore a site
$j_0\in Z^+$ with
\begin{equation}
 \abs{(z_0)_{j_0}}\ge(1+L)^{-C}\sqrt\rho.
 \label{trans:initial-component}
\end{equation}

The other field has the representation
\begin{equation}
 U_0=\frac{A^{-1}p}{\epsilon}
 +A^{-1}M\mathcal B^{-1}
       \frac{M^\top A^{-1}p}{\epsilon}.
 \label{trans:reflected-field}
\end{equation}
The profile estimate and $\abs{\psi_i}\le C\epsilon$ on $Z$ show
$\norm{M^\top A^{-1}p}\le C(1+L)^C\epsilon$. Thus the reflected
term in \Cref{trans:reflected-field} has polynomial norm, with exponent
fixed independently of $D_2$. At depth
$d=\lfloor D_2\log L\rfloor$, the profile has
\[
 \frac{\abs{\psi_i}}\epsilon\ge c e^{\nu d}
 \ge c' L^{\nu D_2}.
\]
Choose $D_2$ large enough that this exceeds both the reflected term
and the profile error in units of $\epsilon$, with fixed slack.
Every site $i$ at that depth then satisfies
\begin{equation}
 \abs{(U_0)_i}\ge2.
 \label{trans:terminal-field}
\end{equation}

\subsection{A path of resolvent entries}

Join $j_0$ to a site at depth $d$ by a physical nearest-neighbor path
of length $O(1+\log L)$ in $J$, moving inward by decreasing absolute
coordinates. The corner clipping guarantees that the path can avoid
$Z$ after its first vertex. For each successive pair $y,j$, with
$j\notin Z$, \Cref{trans:gamma-inverse} gives
\[
 (\Gamma^{-1})_{yj}=1.
\]
We next replace this physical edge by a uniformly bounded number of
large entries of $\Gamma$. This replacement is needed because
$\Gamma$ need not be spatially local.

Write $B_\Gamma=\Gamma^{-1}$, $n_J=\abs J$, and
$K_\Gamma=\max(1,\norm{B_\Gamma})$. Let
\[
 S_y=\supp(B_\Gamma e_y),\qquad
 S_j=\supp(B_\Gamma e_j),\qquad
 \delta_\Gamma=(2n_JK_\Gamma^2)^{-1}.
\]
Form the undirected graph on $S_y\cup S_j$ with an edge between
$a\in S_y$ and $b\in S_j$ whenever
$\abs{\Gamma_{ab}}\ge\delta_\Gamma$. The endpoints $y,j$ belong
to this vertex set because $(B_\Gamma)_{yj}=1$. If they were not
connected, let $F$ be the diagonal indicator of the component of $y$,
extended by zero outside the graph. The inverse commutator identity
would give
\begin{align*}
 1
 &=\abs{[F,B_\Gamma]_{yj}}
   =\abs{(B_\Gamma[F,\Gamma]B_\Gamma)_{yj}}\\
 &\le\delta_\Gamma
       \norm{(B_\Gamma)_{y,:}}_1
       \norm{(B_\Gamma)_{:,j}}_1
 \le\delta_\Gamma n_JK_\Gamma^2=\frac12.
\end{align*}
Indeed every nonzero summand with $F_{aa}\ne F_{bb}$ must cross
components, so its $\Gamma$ entry is below threshold. This
contradiction proves connectivity.

Every graph edge meets $S_j$, which consists of $j$ and at most its
four physical neighbors: $j\notin Z$ makes the nonlocal term in
\Cref{trans:gamma-inverse} vanish in its column. Thus $S_j$ is a
vertex cover of size at most five. A simple connecting path has at
most ten edges, since a simple path with no two successive vertices
outside a set of size five has at most eleven vertices. The other
support $S_y$ lies in the physical star of $y$ together with $Z$.
All these vertices lie in the shell
$L-\abs i_\infty\le D_2\log L+5$; the fixed buffer accommodates
the physical stars.

Concatenate the paths for the physical edges and erase loops. We
obtain distinct sites
\begin{equation}
 j_0,j_1,\ldots,j_s,\qquad
 s\le C(1+\log L),\qquad
 \abs{\Gamma_{j_a j_{a+1}}}\ge(1+L)^{-C},
 \label{trans:resolvent-path}
\end{equation}
all carrying shell jitters, with $j_s$ at depth $d$. Fixed tie rules
make all choices measurable functions of the regular zero-jitter
data alone.

\subsection{Propagation by the independent jitters}

Choose $D_3$ sufficiently large after $D_2$. For every allowed
simultaneous jitter perturbation, its diagonal operator norm is at
most $L^{-D_3}$. Inverse difference applied to $R_0$ therefore
preserves its polynomial norm, the bounds for $\Gamma^{-1}$, and
the path-edge lower bounds, with constant-factor slack. In particular,
if superscripts denote perturbed and zero-jitter values, then
\[
 z_0^{\rm new}-z_0^{\rm old}
 =-\Gamma^{\rm new}\Delta_J z_0^{\rm old},
\]
where $\Delta_J$ is the jitter diagonal.
\Cref{trans:source-size,trans:initial-component} show that the norm
and starting-component bounds also retain constant-factor slack.
The terminal bound retains, say, $\abs{(U_0)_{j_s}}\ge1$ throughout
these perturbations. For this last assertion use
\Cref{trans:reflected-field} and the uniform profile estimate again:
the large total size of $U_0$ is not controlled by an additive
small-perturbation argument.

Condition on all nonpath jitters. Expose the path jitters successively,
setting unexposed ones to zero only in the intermediate calculations.
Let $w=L^{-D_3}$ be their common half-width. At step $a<s$, changing
the diagonal at $j_a$ by $\zeta\in[-w,w]$ gives the exact rank-one
formula
\begin{equation}
 (z_0^{\rm new})_{j_{a+1}}
 =(z_0)_{j_{a+1}}
 -\frac{\Gamma_{j_{a+1}j_a}(z_0)_{j_a}\zeta}
             {1+\Gamma_{j_a j_a}\zeta}.
 \label{trans:jitter-update}
\end{equation}
Values on the right refer to the state just before this update.
The denominator has modulus between $1/2$ and $3/2$ after increasing
$D_3$ if necessary. For $\zeta\in[-w,-w/2]$ and
$\zeta'\in[w/2,w]$, the two values on the right differ by
\[
 \abs{\Gamma_{j_{a+1}j_a}(z_0)_{j_a}}
 \frac{\abs{\zeta-\zeta'}}
 {\abs{1+\Gamma_{j_a j_a}\zeta}
  \abs{1+\Gamma_{j_a j_a}\zeta'}}
 \ge c w(1+L)^{-C}\abs{(z_0)_{j_a}}.
\]
If one of the two intervals contains a value below one quarter of
this separation in modulus, every value in the other interval is
above that threshold. If it contains no such value, that interval
itself is successful. Hence the conditional success probability is
at least $1/4$, and success gives
\begin{equation}
 \abs{(z_0^{\rm new})_{j_{a+1}}}
 \ge(1+L)^{-C}\abs{(z_0)_{j_a}}.
 \label{trans:jitter-gain}
\end{equation}
Constants and the exponent here include the fixed $D_3$.

Finally expose the jitter at $j_s$. Applied to the scalar $F_0$,
the update is
\[
 F_0^{\rm new}=F_0-
 \frac{(U_0)_{j_s}(z_0)_{j_s}\zeta}
      {1+\Gamma_{j_sj_s}\zeta}.
\]
Since $\abs{(U_0)_{j_s}}\ge1$, the same two-interval argument
gives a final success of probability at least $1/4$ and size at least
$(1+L)^{-C}\abs{(z_0)_{j_s}}$. Only the current component is used
at each step; later jitters need not preserve previously propagated
components. There are $O(1+\log L)$ updates, all at distinct sites
and hence using independent uniforms. Combining
\Cref{trans:initial-component,trans:jitter-gain}, we conclude that
\begin{equation}
 \PP_{\rm jitters}\left\{
   \abs{p^\top z_0/\epsilon}
     \ge e^{-C(\log L)^2}\sqrt\rho\right\}
 \ge e^{-C(1+\log L)}.
 \label{trans:jitter-success}
\end{equation}
This holds conditionally on every regular zero-jitter boundary and
$r$ datum, and uniformly after conditioning on nonpath jitters.

\subsection{Restoring absorption and the root}

Reinsert the prescribed artificial absorptions in $J$. If their
diagonal is $\tau A_{\rm art}$, with $0\le A_{\rm art}\le I_J$,
the rootless inverse becomes
\[
 R_\tau=
 \left(R_0^{-1}-i\tau\begin{pmatrix}A_{\rm art}&0\\0&0\end{pmatrix}
 \right)^{-1}.
\]
Write $\Gamma_\tau=(R_\tau)_{JJ}$ and
$z_\tau=(R_\tau)_{JP}\widehat k$. Inverse difference and the
polynomial bounds already established give
\begin{equation}
 \norm{\Gamma_\tau-\Gamma}\le\tau C(1+L)^C,\qquad
 \norm{z_\tau-z_0}\le\tau C(1+L)^C\sqrt\rho.
 \label{trans:absorption-errors}
\end{equation}
Thus the error in $p^\top z_0/\epsilon$ is at most
$\epsilon^{-1}\tau C(1+L)^C\sqrt\rho$.

Now restore the open root. Its scalar Schur denominator is
\[
 d=c(v)-E-i\sigma-i\tau a_0-p^\top\Gamma_\tau p.
\]
For the auxiliary system with internal artificial absorptions omitted,
the calibration and block formula for $\Gamma$ give
\begin{equation}
 \frac{c(v)-E-i\sigma-p^\top\Gamma p}{\sigma}
 =-i-\frac{p^\top A^{-1}M\mathcal B^{-1}M^\top A^{-1}p}{\sigma}.
 \label{trans:root-denominator}
\end{equation}
The boundary profile estimate gives
$\norm{M^\top A^{-1}p}\le C(1+L)^C\epsilon$, and hence the
second term is polynomially bounded. Reinstating absorption changes
the normalized denominator by at most
$\sigma^{-1}\tau C(1+L)^C$, including the absorption at the root.
Because $c_*>2\nu$, this is small for large $L$. We have therefore
proved, uniformly on all the data being used,
\[
 \abs{d/\sigma}\le C(1+L)^C.
\]
Passivity and the open root also give $-\operatorname{Im}d\ge\sigma$,
so this denominator is nonzero.

Block inversion now identifies the amplitude, with an irrelevant sign,
as
\begin{equation}
 t=-\frac{p^\top z_\tau/\epsilon}{d/\sigma}.
 \label{trans:final-root-amplitude}
\end{equation}
The numerator error in \Cref{trans:absorption-errors} is exponentially
smaller than the successful lower bound in \Cref{trans:jitter-success}.
The polynomial denominator and the success probability can therefore
be incorporated into $e^{-C(\log L)^2}$. For each good boundary datum,
integration over the regular $r$ event, of probability at least $1/2$,
and then all jitters gives
\[
 \E_{r,\rm jitters}\abs t^2
 \ge e^{-C(\log L)^2}\rho.
\]
When $\rho=0$, invertibility in \Cref{trans:normalization} gives
$\widehat k=0$, and the inequality is immediate. Finally integrate
over the weighted good boundary event from \Cref{imp:weighted}.
This proves \Cref{trans:lower} and completes the proof of
\Cref{trans:cell}.

For clarity, the choices above have no circular dependence. First fix
the profile and its plateau, the box-width exponents, and $D_1$ as
in \Cref{cell:branch}. The polynomial in \Cref{imp:weighted}, the
regular zero-jitter inverse thresholds, and the exponents for the
profile residual and reflected field can then be fixed without
$D_2,D_3$. Choose $D_2$ to obtain the profile height in
\Cref{trans:terminal-field}. Choose $D_3$ afterwards to guarantee
the profile estimates and all the perturbation margins used above.
The constants in \Cref{trans:jitter-success} may depend on this
last choice. Any fixed $c_*>2\nu$ then makes the absorption errors
small for sufficiently large $L$; no polynomial threshold is chosen
by using the norm of an exterior matrix.

\subsection{An envelope for arbitrary interior choices}

For comparison between different interior laws we need a bound that
does not require the branch on $V$. Keep the same geometry, boundary
data, exterior source, and prescribed artificial mask. Allow any set
$D\subset V$ of internal cell ports, at most one per site, with
strengths $s_d\in[e^{-C_0L},1]$ and real reference potentials bounded
in modulus by a fixed $C_0$. All other internal real diagonals are
their true values, and there are no cell ports on $B$. The set $D$,
strengths, and references may depend arbitrarily on the true
interior potentials and on auxiliary choices.

\begin{proposition}[Interior envelope]\label{trans:envelope}
For almost every fixed boundary datum there is a nonnegative function
$\mathfrak E(v_V)$ depending only on the true interior potentials,
the fixed exterior and boundary data, and the artificial mask, such
that every allowed internal configuration satisfies
\begin{equation}
 \sum_{d\in D}\abs{S_{d,\rm source}}^2
 \le\mathfrak E(v_V),\qquad
 \E_{V,\rm true}\mathfrak E\le e^{C(1+L)^C}\rho.
 \label{trans:envelope-bound}
\end{equation}
The same integrated bound, with adjusted constants, holds under any
density bounded by $e^{C(1+L)^C}$ relative to the true product law on
$V$. The constants are uniform in the exterior data and in all
adaptive interior choices.
\end{proposition}

\begin{proof}
First fix a candidate subset $D$ and let $E_D$ denote its coordinate
inclusion in $V\oplus P$. Replace all its strengths by $1$ and all
its real reference potentials by their true values, retaining the
artificial probes. Write $G_D'$ for the inverse after replacement
and $G_D$ for the inverse of any of the original configurations with
this set of sites. The two inverse matrices differ by
\[
 (G_D')^{-1}=G_D^{-1}+E_D\Delta_D E_D^\top,
 \qquad \norm{\Delta_D}\le C.
\]
Let $s_* = e^{-C_0L}$. Absorption before replacement gives
$\norm{E_D^\top G_D E_D}\le s_*^{-1}$. For any source vector $x$,
inverse difference gives the pointwise identity
\[
 E_D^\top G_Dx
 =\bigl(I_D+E_D^\top G_D E_D\Delta_D\bigr)
      E_D^\top G_D'x.
\]
Consequently, including the factor $4$ for squared scattering
amplitudes and using $s_d\le1$,
\begin{equation}
 \sum_{d\in D}\abs{S_{d,\rm source}}^2
 \le e^{C(1+L)}
       \norm{E_D^\top(G_D')_{VP}\widehat k}^2.
 \label{trans:replacement}
\end{equation}
In this formula and below, $E_D$ is also used for its inclusion in
$V$ when the block dimensions specify this meaning. The replacement
matrix depends on $D$ and true potentials only, so the right side
is independent of the original reference values and strengths.

All these inverses can be justified on a common full-measure set.
For $G_D'$, eliminate $P$ and use the true diagonal densities on $V$;
the real-part determinant is a nonzero diagonal polynomial. If an
original matrix had a kernel vector, absorption would make it vanish
on $D$. Restoring true reference values and replacing strengths at
those sites would leave it a kernel vector of the replacement
matrix. Thus the original inverse exists whenever the replacement
inverse does, even for potential-dependent choices. There are only
finitely many subsets $D$, so their generic sets can be intersected.

Fix $D$. Ward and transpose symmetry imply
\begin{equation}
 X_D:=\norm{(G_D')_{VV}E_D}_{\rm HS}^2
 \le\tr\operatorname{Im}(G_D')_{VV}.
 \label{trans:envelope-ward}
\end{equation}
We bound the expectation of the trace under the true product law
by common-shift phase integration. If $n_V=\abs V$, parameterize
the cube $[-h,h]^{n_V}$ by $v_{n_V}=t$ and
$v_i=t+w_i$ for $i<n_V$. This change has Jacobian $1$; the allowed
$t$ form an interval of length at most $2h$, and each
$w_i\in[-2h,2h]$. Along this interval every true real diagonal in
$V$ shifts at rate $1$. The only other derivatives are
$-i\tau a_i'(v_i)$, which obey
$\abs{a_i'}^2\le Ca_i$. The absorption argument of
\Cref{mat:phase} therefore gives
\[
 \Phi'(t)\ge\frac12\tr\operatorname{Im}(G_D')_{VV}
                -C\dim(V\oplus P).
\]
The phase range and interval length bound the unnormalized trace
integral by $C(1+L)^C$. Integrating over the transverse box and
dividing by $(2h)^{n_V}$ costs at most $e^{C(1+L)^C}$; the constants
may depend on the fixed $h$. Strict regularization followed by
Fatou handles null singular slices. Thus
\begin{equation}
 \E_{V,\rm true}X_D
 \le\E_{V,\rm true}\tr\operatorname{Im}(G_D')_{VV}
 \le e^{C(1+L)^C},
 \label{trans:envelope-phase}
\end{equation}
uniformly in $D$ and in the exterior matrices.

Let $H_V'$ be the $V,V$ block of $(G_D')^{-1}$. It has bounded
operator norm, since its diagonals are true potentials with unit
cell absorption and bounded artificial absorption. The $V$ rows of
$(G_D')^{-1}G_D'=I$ give
\[
 M(G_D')_{PV}E_D
   =E_D-H_V'(G_D')_{VV}E_D.
\]
The lower singular-value bound for $M$ implies the pointwise estimate
\begin{equation}
 \norm{(G_D')_{PV}E_D}_{\rm HS}^2
 \le C\bigl(\abs D+X_D\bigr).
 \label{trans:envelope-boundary-field}
\end{equation}
For the exterior source, use the $V,P$ block of
$G_D'(G_D')^{-1}=I$ and \Cref{trans:normalization}:
\begin{align*}
 E_D^\top(G_D')_{VP}\widehat k
 &=E_D^\top(G_D')_{VP}(\widehat T-iI_P)k_*\\
 &=-E_D^\top(G_D')_{VV}Mk_*
       -iE_D^\top(G_D')_{VP}k_*.
\end{align*}
Transpose symmetry, the bounded norm of $M$, and
\Cref{trans:envelope-boundary-field} now show
\begin{equation}
 \norm{E_D^\top(G_D')_{VP}\widehat k}^2
 \le C\rho\bigl(\abs D+X_D\bigr).
 \label{trans:envelope-source}
\end{equation}
Again, this uses no norm bound on $\widehat T$ or $\widehat k$.
\Cref{trans:envelope-phase,trans:envelope-source} give the
desired integrated bound for each replacement problem.

Finally define the pointwise envelope
\[
 \mathfrak E(v_V)=e^{C(1+L)}
   \sum_{D\subset V}
    \norm{E_D^\top(G_D')_{VP}\widehat k}^2,
\]
with the empty-set summand zero and any value on the common null
exceptional set. There are $2^{\abs V}=e^{O(L^2)}$ summands.
\Cref{trans:replacement} proves pointwise domination for
every adaptive choice, and summing the expectation bounds proves
\Cref{trans:envelope-bound}. Multiplication by a bounded density
gives the final assertion of the proposition.
\end{proof}

When this envelope is compared with the branch lower bound, the
corner-box conditionings must coincide. The single-side laws on
$B$ may differ: their settings and artificial absorptions change
$\rho$ only by the pointwise bounded factors in \Cref{imp:weighted}.
This is the precise boundary comparison needed when an adaptive
collection of descendant ports is replaced by an open parent root.
